\subsection*{A split Birman model for the remaining finite-cover condition}

\statement{B.75}{Proved}{A four-moving-point cocycle has exactly the residual centralizer as its completed zero set}
This is a complementary formulation of the source-membership problem
in the attached LM4 F.14--F.15.  Those results already identify and
solve the full involution cocycles in their specified two-point
interaction kernel; neither that correction nor the formulation below
proves membership in the original source.

Let \(G=C_{\Gamma_{0,[12]}}(J)\cong C_2*\mathbb Z\) be the full
dihedral centralizer of B.74.  In its common-involution quotient use
the ordinary labels
\[
\begin{split}
 U&=\{1,3\},\qquad V=\{2,5\},\\
 (W_1,W_2,W_3,W_4)&=(\{6,8\},\{12,10\},\{9,11\},\{4,7\}).
\end{split}                                                    \tag{B.75.1}
\]
Let \(R_8\) fix \(U,V\) individually, preserve the unordered
branch pair, and permit permutations of the four \(W_i\).
It has index five in B.74's \(\widetilde Q_8\).
Put \(G_V=G\cap R_8\); thus \([G:G_V]=5\) and
\([G_V:H_5]=2\).
Forgetting the four \(W_i\) gives actual split extensions
\[
\begin{split}
 1&\longrightarrow N\longrightarrow R_8
       \xrightarrow{\pi}Q_4\longrightarrow1,\\
 N&=B_4(S_{0,4}),\qquad Q_4=C_2*\mathbb Z,\qquad
 R_8=N\rtimes Q_4,\\
 \widehat R_8&=\widehat N\rtimes\widehat Q_4.
\end{split}                                                    \tag{B.75.2}
\]
Here \(B_4(S_{0,4})\) is the unordered four-strand braid group
on the sphere punctured at the four retained marks, and the
section in (B.75.2) lies in the actual \(G_V\cong C_2*F_3\).

Use the second actual tetrahedral involution
\[
       s=(1\ 2)(3\ 5)(4\ 6)(7\ 8)(9\ 10)(11\ 12).
                                                               \tag{B.75.3}
\]
Its quotient \(\bar s\) normalizes \(R_8\) and acts trivially
on \(Q_4\).  The based nonabelian cocycle
\[
 \delta:G_V\longrightarrow N,\qquad
       \delta(g)=\bar s g\bar s^{-1}g^{-1}
                                                               \tag{B.75.4}
\]
is injective as a map of discrete sets.  Its actual continuous
extension satisfies, under B.74's faithful based embedding,
\[
 \widehat\delta:\widehat G_V\longrightarrow\widehat N,
       \qquad \widehat\delta^{-1}(1)=\mathcal K.
                                                               \tag{B.75.5}
\]
Its fibers are the sets \(h\mathcal K\), namely left cosets;
no normality of \(\mathcal K\) in \(\widehat G_V\) is asserted.

Consequently \(\mathcal K=1\) is equivalent to the following
uniform finite-cover condition:
\[
\begin{gathered}
 \text{for every }U_0\mathrel{\trianglelefteq_f}G_V
 \text{ there is a characteristic }N_0\leq_fN\text{ such that}\\
       g\in G_V,\quad\delta(g)\in N_0
                    \quad\Longrightarrow\quad g\in U_0.
\end{gathered}                                                   \tag{B.75.6}
\]
This condition remains unresolved.

\begin{proof}
\emph{The actual section.}
The four-point forgetful map has the Birman kernel
\(N=B_4(S_{0,4})\): its fiber is the unordered configuration
space of the four forgotten points on the four-punctured
sphere.  There is no exceptional central quotient, since
the retained surface has negative Euler characteristic and
the usual configuration-space Birman sequence applies.

The restriction to \(G_V\) is explicit.  Write \(u\) for
B.74's cyclic source parameter.  In the branch-normalized
coordinate the point \(V\) has value
\(((s+1)/(s-1))^2\), where \(u=s^2\).
The four-point orbifold coordinate is therefore
\[
 w=4\left(\frac{u+1}{u-1}\right)^2.
                                                               \tag{B.75.7}
\]
The orbifold for \(G_V\) is \(\mathbb P^1_w\) with a cone
point of order two at zero and four cusps
\[
       4,\ \infty,\ a=-4-\tfrac85\sqrt5,\
                         b=-4+\tfrac85\sqrt5.
\]
Indeed \(a,b\) are the two images of
\(\mu_5\setminus\{1\}\) under (B.75.7), and the degree-five
map to B.74's \(q\)-orbifold is
\[
 q=\frac{4w(w^2+40w+80)^2}{(w-4)^5},\qquad
 q-4=\frac{16(5w^2+40w+16)^2}{(w-4)^5}.
                                                               \tag{B.75.8}
\]
The actual common basepoint is \(w_*=4/5\), mapping to
\(q_*=-121\).

Forgetting the four \(W_i\) fills exactly \(a,b\).
Choose based cone and cusp meridians, with paths from this
same basepoint, so that the restriction is
\[
 \langle e,a_0,c_1,c_2\mid e^2=1\rangle
       \longrightarrow\langle e,a_0\mid e^2=1\rangle,
       \qquad c_1,c_2\longmapsto1.
                                                               \tag{B.75.9}
\]
Its restriction to \(\langle e,a_0\rangle\) is an
isomorphism.  This provides the claimed section inside
\(G_V\), and proves the discrete splitting in (B.75.2).
It does not retract \(R_8\) onto \(G_V\); the two additional
cusp meridians are killed.

\emph{Full topology of the Birman kernel.}
The group \(N\) is finitely generated.  Its pure subgroup
is an iterated extension of finite-rank free groups by the
Fadell--Neuwirth fibrations, and has finite index in \(N\).
Characteristic finite-index subgroups of \(N\) are cofinal
among its finite-index subgroups.  For such an \(N_0\), the
\(Q_4\)-action preserves \(N_0\) and factors on \(N/N_0\)
through a finite automorphism group.  The resulting finite
semidirect quotient of \(R_8=N\rtimes Q_4\) induces exactly
the quotient \(N/N_0\) on \(N\).
Thus \(N\) has its full profinite topology in \(R_8\).
The retraction and section give the full topology on the
chosen \(Q_4\) as well.  Right exactness of completion for
the normal quotient, together with the embedded \(\widehat N\),
proves the last equality in (B.75.2).
This argument does not complete an infinitely generated
abstract source kernel in place of its embedded closure.

\emph{The exact involution and its deck discrepancy.}
The actual \(s\) belongs to \(T\), commutes with \(r\), and
generates \(A_5\) together with \(J\).
Its quotient swaps \(U,V\), swaps \(W_1,W_4\) and
\(W_2,W_3\), and swaps the branch pair.  The last assertion
also follows from the fact that distinct commuting sphere
involutions exchange each other's two fixed points.
Thus \(\bar s\) normalizes \(R_8\).
On the four retained points it is the double transposition
\((U\ V)(\text{branch}_1\ \text{branch}_2)\).
The pure four-point group centralizes the canonical Klein
four subgroup.  The additional allowed branch transposition
also fixes this double transposition.  Hence conjugation
by \(\bar s\) acts trivially on \(Q_4\), proving that
(B.75.4) lands in \(N\).  Direct multiplication gives
\[
       \delta(gh)=\delta(g)\,g\delta(h)g^{-1}.
                                                               \tag{B.75.10}
\]
Conjugation is an actual automorphism of \(R_8\), so the
same construction extends continuously into the actual
completed kernel \(\widehat N\).

The full preimage of \(R_8\) in the symmetric twelve-point
group is \(\langle r\rangle\times R_8\), using the unique
lift fixing original label one.  B.73's prime-order
comparison gives the same exact finite deck kernel after
completion.  The element \(s\) normalizes this preimage,
but need not preserve the chosen section.  Thus a zero
of \(\widehat\delta\) initially implies only that its
actual lift conjugates \(s\) to \(s\) or \(rs\).
The second possibility is excluded by the exact finite
orders in the original action:
\[
                      |sd|=3,\qquad |rsd|=5.
                                                               \tag{B.75.11}
\]
Every element of \(\widehat G_V\) centralizes \(d\).
It cannot therefore conjugate \(s\) to \(rs\), since this
would conjugate elements of different orders in the actual
finite \(A_5\).  This covers the nonpure coset of
\(H_5\) in \(G_V\), whose permutation swaps original
labels two and five; purity has not been silently assumed.

\emph{Both zero-set comparisons.}
A discrete zero now centralizes \(\langle d,r,s\rangle=A_5\),
whose discrete centralizer is trivial by B.52.
Moreover \(\delta(g)=\delta(h)\) is equivalent to
\(h^{-1}g\) being fixed by conjugation by \(\bar s\).
The same deck argument proves the stated discrete
set-injectivity.

A completed zero likewise centralizes the actual \(A_5\).
The permutation centralizer of its transitive twelve-point
action has order two, and its nonidentity swaps original
labels one and three.  But every element of
\(\widehat G_V\) fixes these labels individually.
The zero is therefore pure and belongs to
\(C_{\widehat P}(A_5)=\mathcal K\).
Conversely \(\mathcal K\subseteq\widehat H_5\subseteq
\widehat G_V\) and centralizes \(s\), so its specified
based quotient has cocycle value one.  This proves
(B.75.5), and the preceding equality-of-values calculation
gives its precise left-coset fibers.

Finally, if the completed zero fiber is trivial, compactness
of the complement of \(\widehat U_0\) supplies an open
identity neighborhood in \(\widehat N\) disjoint from its
image.  Cofinality of the characteristic finite-index
subgroups gives (B.75.6).
Conversely each such finite condition passes from the dense
discrete source to the completed source because both sets
are clopen.  A completed zero then belongs to every
\(\widehat U_0\), and hence is one.
Pointwise residual-finiteness separation of individual
nonidentity values \(\delta(g)\) does not establish this
uniform condition.
\end{proof}

The finite and rational inputs are freshly checked from the
original marked permutation data by
\path{work_DE/next_common_involution/check_Birman_cocycle_data.py};
its actual output and complete certificate are retained beside it.
They verify (B.75.11), the full permutation centralizer and its
pole swap, the nonpure \(V\)-stabilizer, the rational maps,
and the common basepoint.  The completed comparison is the
proof above, not an inference from those finite tests.

